% Conclusion. Hypotheses as in CLAUDE.md §11.4; every verdict points to the table that holds the
% numbers (rule 11). Verdict words are the scripts' (experiments/results/*/..._verdicts.json,
% e13_summary.json, bench/results/pipeline_bench.json) or "not run".

\chapter{Conclusion}
\label{ch:conclusion}

\section{Evidence for each hypothesis}
\begin{table}[h]\centering\small
\begin{tabular}{lp{0.42\linewidth}p{0.36\linewidth}}
\hline
& Hypothesis & Outcome \\
\hline
H1 & Audited ACI (G5) keeps accepted large errors below target for every proposer; the baselines
exceed it as proposers become aggressive &
G5 never exceeded $\alpha$ in any run at $\tau = 0.2$~K (Table~\ref{tab:e7-tau02}); the baselines
did not fail as predicted (only G4 did, through selection bias); at the pre-registered $\tau = 1$~K
the test could not separate the gates (Table~\ref{tab:e7-tau1}). Not supported as stated. \\
H2 & At equal exact-solve budget G5 finds designs no worse than G1 & not run (E8) \\
H3 & H1 replicates on Darcy flow & not run (E9) \\
H4 & The calibrated signal lowers designers' accepted-wrong rate & not run (no ethics approval) \\
H5 & The GPU pipeline certifies at least ten times more designs per second &
passed (Table~\ref{tab:e4}) \\
H6 & Certified early stopping saves at least 50\,\% of PCG iterations at 0.1~K &
passed with the improved pointwise certificate (Table~\ref{tab:e13}) \\
H7 & Certified audits keep validity with at least 50\,\% fewer audit iterations &
validity kept, saving just under 50\,\%: not supported (Table~\ref{tab:e14-tau02}) \\
H8 & The per-group gate G6 holds per group where G5 does not & not supported
(Table~\ref{tab:e15-tau02}) \\
H9 & Risk-adaptive audits keep validity with 30\,\% fewer audits & not supported
(Table~\ref{tab:e16-tau02}) \\
H10 & The verified solver predicts measured boards & dropped with the hardware validation \\
\hline
\end{tabular}
\caption{Hypotheses of the project plan and what the experiments show.}
\label{tab:hypotheses}
\end{table}

\section{What the project shows}
The certificate side of TwinHeat works as intended on the real four-layer problem. Every test board
received a certificate, no certified bound was ever below the true error, and stopping when the
error is certified rather than at a fixed tolerance saves more than half of the solver's work at a
0.1~K guarantee. One batched GPU launch certifies designs more than thirty times as fast as the
one-at-a-time CPU loop. Getting there required two changes to the planned construction, both
justified by the lemma itself rather than by tuning: the certificate's auxiliary vectors had to be
computed by multigrid-preconditioned conjugate gradients instead of a few multigrid cycles, and
any such vector is valid as long as the positivity conditions are checked before use.

The statistical side holds where it was tested: the audited gate kept the accepted-large-error
rate below its target in every run once the threshold made large errors common, and the
flagged-only ablation failed for the reason the theory predicts. But the experiments did not show
the separation from the baselines that the plan expected, the high-probability bound is vacuous at
our run lengths, and at the threshold where the gates differ they verify most designs, because the
label-free uncertainty scale is not sharp enough to accept many of them. The per-group and
risk-adaptive extensions did not deliver.

\section{Future work}
\begin{itemize}
\item Run the certified verification of the live server on the GPU path, which the batched
benchmark shows is fast enough to approach the interactive latency targets.
\item A sharper label-free scale, so that the gate can accept more designs at a tight threshold.
\item The per-group guarantee: a proof for the MultiValid gate, or a different construction with
enough audits per group.
\item The efficiency comparison at equal budget with feedback to the proposer (E8), the Darcy
replication (E9), and the user study once ethics approval is obtained.
\item Viewer: flux particles and the remaining visual encodings of uncertainty.
\end{itemize}
